\section*{Budget Justification}

\textit{Draft basis. The amounts below match the companion formula-driven
workbook and total \$486,436 over three years. Salary, fringe, and F\&A values
are planning assumptions and must be replaced with Georgia Tech's approved
rates before submission. No voluntary committed cost sharing is proposed.}

\subsection*{A. Senior Personnel}

\textbf{PI Brian Kyanjo:} 7.2 person-months (60\% of a 12-month appointment)
in each project year. The PI will direct the project and lead model-adapter and
PISM development, ICESEE/CryoBench integration, secure execution, releases,
operations, and technical documentation. The workbook assumes a Year-1 base
salary of \$85,000 with 3\% annual escalation, yielding \$51,000, \$52,530,
and \$54,106 in requested salary.

\textbf{Co-PI Alexander A. Robel:} one summer month in each project year. The
Co-PI will lead science requirements and validation, the observational and
cross-model demonstrations, user evaluation, and community engagement. The
workbook assumes a nine-month academic salary of \$130,000 with 3\% annual
escalation, yielding \$14,444, \$14,878, and \$15,324.

Total senior-personnel salary is \$202,282 over three years. Final person-months
and institutional base salaries will be certified by Georgia Tech.

\subsection*{B. Other Personnel and C. Fringe Benefits}

No postdoctoral scholar, student, or separate research software engineer is
included in the current planning budget. The PI provides the principal
software-engineering effort. Fringe benefits are estimated at 30\% of salary,
or \$60,685 over three years. The final budget will use the applicable Georgia
Tech fringe rates by appointment type. If a CI-professional position is added,
the budget and required CI Professional Development Plan will be revised
together.

\subsection*{D. Equipment}

No equipment is requested. Computing will use existing Georgia Tech PACE and
web infrastructure, user-managed resources where appropriate, and budgeted
cloud services.

\subsection*{E. Travel}

\$14,400 is requested across three years (\$4,800 per year). Each year includes
one trip near NSF for the required CSSI PI meeting, estimated at \$1,800, and
\$3,000 for presentation, workshop delivery, or stakeholder engagement at a
cryosphere, PISM, or science-gateway venue. Costs cover economy airfare or
ground transportation, lodging, per diem, and registration in accordance with
Georgia Tech travel policy.

\subsection*{F. Participant Support Costs}

\$18,000 is requested (\$6,000 per annual workshop) for travel, lodging,
registration, or stipends for external participants. Selection will prioritize
students, early-career researchers, minority-serving institutions, EPSCoR
jurisdictions, and institutions without dedicated HPC/software staff. These
funds are excluded from the MTDC base and will not be rebudgeted to other
categories without required NSF approval.

\subsection*{G. Other Direct Costs}

\$21,000 is requested (\$7,000 per year): \$1,500 for materials and supplies
supporting workshops, software testing, and documentation; \$1,500 for
open-access publication and archival dissemination; and \$4,000 for bounded
AWS compute/storage/network charges used to qualify and demonstrate the cloud
backend. Cloud usage will be protected by quotas, budgets, and alerts. The
bounded computing-access pilot described in the Project Description (planning
scenario of about \$1,560: 50 participants at \$25 with 25\% contingency) is
not included in these totals; it requires institutional approval and cost
classification before it is added. Development of the allocation service
(Task 2.4) is part of the PI effort above; administering allocations and
reconciling charges add operational effort that is not separately budgeted. Final
classification of publication and cloud charges will follow institutional and
NSF policy.

\subsection*{I. Indirect Costs and Cap Check}

The planning workbook applies a provisional 57\% F\&A rate to Modified Total
Direct Costs, excluding participant support, producing estimated indirect
costs of \$170,069. The final submission will use Georgia Tech's current
federally negotiated rate and correct base. Estimated annual totals are
\$158,098, \$162,105, and \$166,233; the three-year total of \$486,436 is
\$113,564 below the \$600,000 Elements cap. Any personnel, subaward, rate, or
scope change will be reconciled across the Research.gov forms, workbook,
budget justification, management plan, and CI-professional plan.
